\section{Граничные условия}

Для замыкания дискретной задачи необходимо задать граничные условия на давление~$p$ и (при использовании модели JFO) на степень заполнения~$\vartheta$. Тип граничных условий определяется физической постановкой (подразделы~2.1.2--2.1.4) и геометрией расчётной области. Ниже систематизированы три основных типа граничных условий для давления.

\textit{Дирихле (заданное давление).}
На соответствующем участке границы $\Gamma_D$ давление фиксируется:
\[
  p = p_0 \quad \text{на} \quad \Gamma_D.
\]
Наиболее распространённый случай -- $p_0 = p_{\mathrm{cav}} = 0$ (в избыточных давлениях) на открытых границах расчётной области, где смазочный слой сообщается с окружающей средой. Типичные примеры: торцы по оси~$z$ для радиального подшипника, радиальные границы $r = R_{\mathrm{in}},\, R_{\mathrm{out}}$ для упорного.

\textit{Периодичность.}
По окружному (периодическому) направлению давление совпадает на противоположных границах области:
\[
  p(x_{1,s},\, x_2) = p(x_{1,t},\, x_2).
\]
Данное условие применяется в упорных и радиальных подшипниках, где расчётная область охватывает полный период по окружной координате. При расчёте сектора вместо периодичности на границах сектора задаются условия симметрии ($\partial p / \partial n = 0$) или заданное давление в зависимости от постановки.

\textit{Непроницаемость (Нейман).}
На участках с условием симметрии или непроницаемости нормальная производная давления обращается в нуль:
\[
  \frac{\partial p}{\partial n} = 0 \quad \text{на} \quad \Gamma_N.
\]
Условие $\partial p / \partial n = 0$ эквивалентно отсутствию напорного потока через границу и используется на линиях симметрии или на искусственных границах при расчёте сектора.

В таблице~\ref{tab:bc_summary} сведены граничные условия для давления, соответствующие трём постановкам, сформулированным в подразделах~2.1.2--2.1.4.

\begin{table}[!ht]
\centering
\caption{Граничные условия для давления в различных постановках
  (координаты $x_1$,~$x_2$ определены в подразделах~2.1.6 и~2.2.1)}
\label{tab:bc_summary}
\begin{tabular}{|l|c|c|}
\hline
\multicolumn{1}{|c|}{\textbf{Постановка}} & \textbf{По $x_1$}
                       & \textbf{По $x_2$} \\
\hline\hline
Упорная
  & периодичность
  & $p = p_{\mathrm{cav}}$, $x_2 = r$ \\
\hline
Радиальная
  & периодичность
  & $p = p_{\mathrm{cav}}$, $x_2 = z$ \\
\hline
Возвр.-поступат.
  & $p = p_{\mathrm{cav}}$, $x_1 = z$
  & периодичность \\
\hline
\end{tabular}

\medskip
\small
Здесь $x_1 = R\theta$ -- окружная координата;
$x_2$ -- вторая координата ($r$ или $z$ в зависимости от постановки).
\end{table}

При использовании массо-сохраняющей модели кавитации JFO (раздел~2.3) дополнительно задаются граничные условия на степень заполнения~$\vartheta$. На границах подачи смазки полагается $\vartheta = 1$ (полное заполнение зазора). По периодическим направлениям степень заполнения периодична, аналогично давлению. На выходных границах с $p = p_{\mathrm{cav}}$ величина~$\vartheta$ явно не фиксируется; она определяется из уравнения сохранения~\eqref{eq:jfo_conservation} и схемы переноса, что соответствует свободному выходу смеси.

Во всех рассматриваемых постановках кавитационное давление принимается равным нулю в избыточных давлениях: $p_{\mathrm{cav}} = 0$.
